\documentclass[10pt,a4paper]{amsart}
\usepackage[margin=19mm]{geometry}
\usepackage{amsmath,amssymb,mathtools}
\usepackage[scr=rsfs,cal=euler]{mathalfa}
\usepackage{xfrac}
\usepackage[unicode,hidelinks,bookmarks=false]{hyperref}
\newcommand{\ie}{i.e.\ }
\newcommand{\cf}{cf.\ }
\newcommand{\resp}{respectively\ }
\newcommand{\Subsection}{\subsection}
\providecommand{\nobreakdash}{\nobreak\hbox{-}}
\setlength{\emergencystretch}{2em}
\multlinegap0pt
\usepackage{bm}
\usepackage{bbm}
\usepackage{dsfont}
\usepackage{xspace}
\usepackage{cancel}
\usepackage{mathtools}
\usepackage{amsthm}
\makeatletter
\@ifundefined{thm}{\newtheorem{thm}{Thm}}{}
\@ifundefined{lem}{\newtheorem{lem}[thm]{Lemma}}{}
\makeatother
\makeatletter
\newcommand{\F}{\mathbb{F}}
\expandafter\ifx\csname poc\endcsname\relax
\newenvironment{poc}{\begin{proof}[Proof of claim]\renewcommand*{\qedsymbol}{$\blacksquare$}}{\end{proof}}
\fi
\makeatother
\title[Intersecting families and nonvanishing multivariate polynomi]{Intersecting families and nonvanishing multivariate polynomials over finite fields}
\author{Shamil Asgarli, Bence Csajb\'{o}k, Chi Hoi Yip}
\date{Source: arXiv:2608.17785v1 (CC BY 4.0)}
\begin{document}
\maketitle

\noindent\emph{Source and licence.} Intersecting families and nonvanishing multivariate polynomials over finite fields. Version arXiv:2608.17785v1 (CC BY 4.0), distributed under the Creative Commons Attribution 4.0 International licence, as recorded in the arXiv licence metadata for this exact version and shown on the abstract page https://arxiv.org/abs/2608.17785v1. Full rights notice: licenses/arxiv-2608.17785-CC-BY-4.0.txt.\par\smallskip

\noindent\textit{Editorial scope.} Counted unit: the Lem whose source label is \texttt{lem:parity} in the article (source0.tex lines 333--339, source label \texttt{lem:parity}), delivered together with the article's own complete original proof of it (source0.tex lines 341--391, one \texttt{proof} environment of 51 lines). No mathematics is written, repaired or re-derived: statement and proof are the article's own text line for line. Required context retained uncounted: none. Every definition, display and label of those lines is retained unchanged. Omitted from this package and not redistributed: 1--332, 340--340, 392--1448. Nothing inside a retained excerpt is omitted or altered: every retained body is delivered exactly as the article's lines are written, so no omission receipt of any kind accompanies this package. Citations: the retained excerpts carry no citation command at all, so no bibliography entry of the article is transcribed and no reference is redistributed. Reference adaptation: none was needed and none was made; every reference inside the delivered unit resolves inside it, so no unresolved reference or citation is produced. Printed slips kept literal: no printed word, symbol, hypothesis or number of the counted statement or of its proof was altered. Layout and class: the article's own class is \texttt{amsart} with packages including fontenc, inputenc, babel, fullpage, times, colonequals, amsmath, amssymb, amsthm, mathtools, url, xcolor, comment, bbm; the wrapper is the lane's pinned \texttt{amsart} editorial wrapper at 10pt on A4, which restates the theorem-like environments the retained text needs and adds the article's own macro definitions that the retained text needs (\texttt{\textbackslash F}), with extra wrapper packages (bm, bbm, dsfont, xspace, cancel, mathtools). The delivered numbering is the unit's own. Assets: no retained span contains a figure environment, a graphics-inclusion command, a PDF-page inclusion, a table or a TikZ picture; no image file is copied into this package, the package declares no asset dependency and no image file is redistributed with it. Licence: version arXiv:2608.17785v1 of this article is distributed under the Creative Commons Attribution 4.0 International licence, as recorded in the arXiv licence metadata for this exact version and shown on the abstract page https://arxiv.org/abs/2608.17785v1. A full rights notice is delivered at licenses/arxiv-2608.17785-CC-BY-4.0.txt. Characters: every retained excerpt is pure ASCII, so the delivered engine, XeLaTeX with its default Latin Modern text font, reports no missing character.
\begin{lem}\label{lem:parity}
Let $q=2^r$, write $\mathbf x=(x_1,\ldots,x_d)$, and let $f\in \F_q[x_1,\ldots,x_d]$ with $\deg f\leq d$. Put $a=[x_1x_2\cdots x_d]f$ and let
\[
  Z(f)=\bigl|\{\mathbf x\in \F_q^d:f(\mathbf x)=0\}\bigr|.
\]
Then $Z(f)\equiv a^{q-1}\pmod2$. In particular, if $a\neq0$, then $f$ has a zero in $\F_q^d$.
\end{lem}

\begin{proof}
Let $\overline{Z}(f)$ be the reduction of $Z(f)$ mod $2$. We view $\overline{Z}(f)$ as an element of $\F_2\subseteq\F_q$. Observe that
\[
\overline{Z}(f)=\sum_{\mathbf x\in\mathbb F_q^d}\bigl(1-f(\mathbf x)^{q-1}\bigr).
\]
Moreover,
\[
\sum_{t\in\mathbb F_q}t^m=
\begin{cases}
1,& m>0\ \text{and}\ q-1\mid m,\\
0,&\text{otherwise}.
\end{cases}
\]
Consequently, $\sum_{\mathbf{x}\in\F_q^{d}} \mathbf{x}^{\boldsymbol{\alpha}}=0$ if $\alpha_i<q-1$ for some $i$. As $\deg f^{q-1}\le d(q-1)$, the only monomial of $f^{q-1}$ that
survives after summing over $\mathbb F_q^d$ is
$x_1^{q-1}\cdots x_d^{q-1}$. We obtain
\begin{equation}\label{eq:Z(f)}
\overline{Z}(f)=[x_1^{q-1}\cdots x_d^{q-1}]f^{q-1}.
\end{equation}
It remains to compute this coefficient. Write $f=\sum_{\boldsymbol{\beta}} c_{\boldsymbol{\beta}}\mathbf{x}^{\boldsymbol{\beta}}$ so that $c_{\mathbf 1}=a$, where $\mathbf{1}=(1,\ldots,1)$. A term in the multinomial expansion that
contributes to $\mathbf x^{(q-1)\mathbf1}$ is specified by
nonnegative integers $m_{\boldsymbol\beta}$ satisfying
\[
\sum_{\boldsymbol{\beta}} m_{\boldsymbol{\beta}}=q-1,
\qquad
\sum_{\boldsymbol{\beta}} m_{\boldsymbol{\beta}}\boldsymbol{\beta}=(q-1)\mathbf 1.
\]
Taking total degrees shows that every $\boldsymbol{\beta}$ with $m_{\boldsymbol{\beta}}>0$ satisfies $|\boldsymbol{\beta}|=d$.

Suppose the corresponding multinomial coefficient $\binom{q-1}{(m_{\boldsymbol{\beta}})_{\boldsymbol{\beta}}}$ is odd. By Kummer's theorem, the 
addition $\sum_{\boldsymbol{\beta}}m_{\boldsymbol{\beta}}=q-1$ is carry-free in base $2$. Since \(q-1\) is odd,
exactly one \(m_{\boldsymbol{\beta}}\) is odd. Reducing the second equality modulo \(2\)
shows that its exponent vector satisfies
$\boldsymbol{\beta}\equiv\mathbf 1\pmod 2$. 
But \(|\boldsymbol{\beta}|=d\), so its \(d\) coordinates are positive odd integers with
sum \(d\). Hence \(\boldsymbol{\beta}=\mathbf 1\).

Replace $m_{\mathbf 1}$ by $m_{\mathbf 1}-1$, and then divide all the $m_{\boldsymbol{\beta}}$ values by $2$. The same argument then applies with $q-1$ replaced by $(q-2)/2=2^{r-1}-1$. Repeating this for all binary digits of $q-1$ shows that
\[
m_{\mathbf 1}=q-1, \qquad \text{and}
\qquad
m_{\boldsymbol{\beta}}=0 \quad \text{for all} \quad \boldsymbol{\beta}\ne\mathbf 1.
\]
Therefore, every other multinomial coefficient is even, and
\begin{equation}\label{eq:coeff-c1}
[x_1^{q-1}\cdots x_d^{q-1}]f^{q-1}
=c_{\mathbf 1}^{q-1}=a^{q-1}.
\end{equation}
Combining equations~\eqref{eq:Z(f)} and \eqref{eq:coeff-c1} yields $\overline{Z}(f)=a^{q-1}$, 
as required.
\end{proof}

\end{document}
